% =====================================================================
% sec-abstract.tex — 摘要与关键词
% =====================================================================
\phantomsection
\addcontentsline{toc}{section}{摘要}
\begin{center}
  {\heiti\zihao{3} 摘\quad 要}
\end{center}
\vspace{0.5em}

现代深度分类神经网络在医疗诊断与自动驾驶等高风险决策场景中广泛应用，然而其预测概率普遍存在严重的「过度自信」现象，导致模型输出的高置信度无法真实反映预测正确的实际概率，削弱了决策系统的统计可靠性。为对这一概率失配问题进行严密的数理统计推断，本文将多分类 Softmax 后处理校准方法 Temperature Scaling 形式化纳入广义线性模型与极大似然估计（MLE）框架，旨在检验深度网络是否存在统计显著的校准偏差并构建完整的推断与检验体系。在理论层面，本文推导了逆温度参数 $\beta = 1/T$ 的对数似然函数、解析 Score 函数与观测 Fisher 信息量，通过条件方差非负性证明了似然函数的严格凹性与参数解的全局唯一性，并采用带回溯线搜索的阻尼 Newton--Raphson 迭代法进行数值求解；进一步借助一阶 Delta 方法推导了温度估计 $\hat T$ 的渐近方差并构造渐近置信区间，辅以非参数 Bootstrap 进行有限样本精度对照；构建了针对无需校准原假设 $H_0: T=1$ 的 Wald、Score 与似然比（LRT）三大渐近检验，并引入双参数 Platt Scaling 作为对照模型。在实证层面，基于 CIFAR-10 官方测试集（$N=10{,}000$ 样本，按 $2{,}000:8{,}000$ 分层划分为校准集与独立评估集，随机种子固定为 20261111），ResNet-20 与 ResNet-56 的温度估计分别为 $\hat T = 1.6245$（95\% Delta CI $[1.5147, 1.7343]$）与 $\hat T = 1.6904$（$[1.5884, 1.7925]$）；三大检验统计量均在 $p < 10^{-16}$ 显著性水平下拒绝原假设；在独立评估集上，两模型的期望校准误差（ECE）分别下降 $66.1\%$ 与 $68.6\%$，对数损失（NLL）分别降低 $17.7\%$ 与 $24.7\%$，Brier 分数显著改善（配对 Wilcoxon 检验 $p < 10^{-16}$，配对 $t$ 检验 $p < 10^{-14}$），且 McNemar 检验证实 Top-1 分类准确率保持不变（$p = 1.000$）。实证结果确证了极大似然估计满足理论 $O(n^{-1/2})$ 渐近相合速率，阐明了三大检验统计量经验排序的曲率机理，表明单参数 Temperature Scaling 在保持判别精度的前提下兼具良好的统计推断效能与实用性。

\keywords{置信度校准；极大似然估计；似然比检验；Delta 方法；Temperature Scaling}
